\documentclass[11pt,a4paper]{article}
\usepackage[margin=25mm]{geometry}
\usepackage{amsmath,amssymb,amsthm,hyperref}
\setlength{\emergencystretch}{2em}
\newtheorem{lemma}{Lemma}
\newtheorem{theorem}[lemma]{Theorem}
\newtheorem{corollary}[lemma]{Corollary}
\newcommand{\gen}[1]{\langle#1\rangle}
\title{Counting generators by their actual orbit requirements}
\author{Human Compression project: a component of the L019 recovery}
\date{}
\begin{document}\maketitle
Let $G=\gen{\rho_0,\ldots,\rho_{r-1}}$ be a transitive
permutation group on $\Omega$, with every maximal omission
$G_j=\gen{\rho_l:l\ne j}$ intransitive. The generators are
involutions. We will impose the string relations where they
are needed. Suppose colour $i$ has exactly one transposition
$\{a,b\}$ joining different $G_i$-orbits. The acquired unique
crossing lemma gives exactly two such orbits $A,B$, containing
$a,b$ respectively. Every $\rho_j$ with $|j-i|>1$ fixes
$a,b$. All actions and orbit sets below are actual permutation
actions; projected intersection properties are not assumed.

\begin{lemma}[Minimal transitivity]\label{minimal}
Suppose a specified family of involutions has exactly $c$
orbits on a set of $d$ points. It contains a subfamily of at
most $d-c$ members with the same orbits. If a subfamily of the
marks of $G_i$ is transitive on each of $A,B$, it contains
every mark of $G_i$.
\end{lemma}
\begin{proof}
Choose an inclusion-minimal subfamily with the same orbit
partition. Deleting any member $s$ splits one of its orbits;
choose one transposition of $s$ between two resulting orbits.
The chosen edges form a forest, with parallel pairs counted
as cycles. Indeed a cycle containing
the chosen $s$-edge would connect its endpoints using only
other chosen edges, hence using elements of the family with
$s$ deleted. This is impossible. The forest has at least
$c$ components and therefore at most $d-c$ edges.

For the second assertion, if $\rho_j$ were absent from the
subfamily, then $G_j$ would contain that subfamily and
$\rho_i$. The subfamily is transitive on both sides, and
the crossing transposition joins them. Thus $G_j$ would be
transitive on $\Omega$, a contradiction.
\end{proof}

\begin{theorem}[Two imprimitive sides]\label{both}
If the actions of $G_i$ on $A,B$ preserve respectively
$m_1,m_2\ge2$ blocks of sizes $k_1,k_2\ge2$, then
\[
                    r\le m_1+m_2+k_1+k_2-3\le(n+2)/2,
       \qquad n=|\Omega|.
\]
This assertion requires neither the string relations nor IP.
\end{theorem}
\begin{proof}
Choose a minimal subfamily $M$ of the omission marks which
is transitive on each of the two sets of blocks. Apply the
first part of Lemma~\ref{minimal} to the simultaneous block
action, which has exactly two orbits. It gives
$|M|\le m_1+m_2-2$. This is minimality for transitivity,
not a claim of independence or IP in the block image.
Each point orbit of $\gen M$ meets every block in the
corresponding side. Consequently the numbers $c_1,c_2$ of
these point orbits satisfy $c_j\le k_j$.

For each remaining omission mark $\rho_l$, choose an actual
edge crossing two $G_l$-orbits. Such an edge exists because
$G_l$ is intransitive but $G$ is transitive. Contract the
$\gen M$-orbits. Its endpoints remain distinct, since
$\gen M\le G_l$. All the selected edges stay within $A$
or within $B$. They form a forest: an alternative path
round a cycle lifts through the $\gen M$-orbits to a path
using no mark of colour $l$, contrary to its eligibility.
Thus the $r-1-|M|$ remaining marks give at most
$c_1+c_2-2\le k_1+k_2-2$ edges. This proves the first bound.
For the second, expand
\[
 n+2-2(m_1+m_2+k_1+k_2-3)
              =(m_1-2)(k_1-2)+(m_2-2)(k_2-2)\ge0.
\]
\end{proof}

\begin{lemma}[Two commuting factors]\label{factors}
Suppose the omission marks are partitioned into two nonempty
subfamilies of sizes $p,q$, generating $H,K$ respectively.
Suppose $H$ and $K$ centralize one another and their product is transitive
on $B$. Assume either $A$ is a singleton or each of $H,K$
is transitive on $A$. Let $h$ be the common size of the
$H$-orbits in $B$ and let $t$ be their number. Then
\[
       p\le h-1,\qquad q\le t-1,\qquad
       h,t\ge2,\qquad p+q\le |B|/2.
\]
\end{lemma}
\begin{proof}
Because $H$ and $K$ commute, $K$ permutes the $H$-orbits
equivariantly. Their product is transitive on $B$, so $K$
is transitive on this orbit set and transports each actual
$H$-action to every other one. In particular all these orbits
have size $h$, and $|B|=ht$.

Choose from the specified $H$ marks a subfamily $M_H$ that
is transitive on one $H$-orbit, with $|M_H|\le h-1$ by
Lemma~\ref{minimal}. Equivariance makes it transitive on
every $H$-orbit. Thus $\gen{M_H,K}$ is transitive on $B$.
It is transitive on $A$ too: $A$ is either a singleton or
$K$ already acts transitively there. The second part of
Lemma~\ref{minimal} forces every $H$ mark into $M_H$,
proving $p\le h-1$.

Choose a subfamily $M_K$ of the specified $K$ marks
transitive on the $t$ actual $H$-orbits, with
$|M_K|\le t-1$. Then $\gen{H,M_K}$ is transitive on $B$
and on $A$, using $H$ on $A$ when needed. Again minimality
forces every $K$ mark into $M_K$, so $q\le t-1$.
Since $p,q$ are positive, $h,t\ge2$. Finally
\[
 p+q\le h+t-2\le ht/2=|B|/2,
\]
where the middle inequality is $(h-2)(t-2)\ge0$.
No assertion of faithful factor restriction or quotient IP
has entered the argument.
\end{proof}

\begin{corollary}[A singleton and the two-point exception]\label{small}
Assume the string relations and $0<i<r-1$.
If $|A|=1$, then $r\le(n+1)/2$.
If $|A|=2$ and both $\rho_{i-1},\rho_{i+1}$ act nontrivially
on $A$, then $r\le n/2$.
\end{corollary}
\begin{proof}
Take $H=G_{<i}$ and $K=G_{>i}$. They commute and have
$p=i$, $q=r-i-1$ specified marks. In the first case apply
Lemma~\ref{factors} with singleton $A$ and $|B|=n-1$.
In the second both factors contain the transposition on
$A$, so each is transitive there; apply the same lemma
with $|B|=n-2$. In each case $r=p+q+1$ gives the claim.
\end{proof}

\begin{theorem}[An interval side in the strict range]\label{orientation}
Assume the string relations, a connected noncommutation
diagram, and $2r>n+2$. At a unique crossing, either the split
is perfect after interchanging its two sides, or one side
is primitive of degree at least two and its nontrivial
distinguished restrictions occupy one interval adjacent
to $i$. After a possible reversal, that side is $B$ and
$G_{<i}$ fixes $B$ pointwise. If the split is not perfect,
some label greater than $i$ acts nontrivially on both sides.
\end{theorem}
\begin{proof}
Both sides cannot be imprimitive by Theorem~\ref{both}.
A singleton side at an interior crossing is excluded by
Corollary~\ref{small}; at an endpoint crossing it immediately
gives a perfect split, with the empty side of the index set
assigned to the singleton. Thus we may choose a nontrivial
primitive side, called $B$.

Suppose its active index set is disconnected. Partition it
at a gap. The two generated images commute and are nontrivial
normal subgroups of their primitive product. Each is therefore
transitive. A nonidentity element of either factor centralizes
the other transitive factor and has no fixed point. Hence
every active distinguished restriction moves $b$. The unique
crossing lemma confines its index to $i-1$ or $i+1$.
There are exactly two active marks, one in each factor,
and each alone is transitive. Thus $|B|=2$ and both neighbours
are active, contradicting Corollary~\ref{small} with the
sides interchanged. Therefore the active set is one interval.
It avoids $i$ and contains a mark moving $b$, so it is
adjacent to $i$. Reverse all indices if necessary to put
it on the right; then $G_{<i}$ fixes $B$ pointwise.

Suppose no mark is shared by the two sides. Every omission
mark is nonidentity on exactly one side. On either interval
of indices separated by $i$, this side cannot change at
consecutive positions: their supports would be disjoint,
so those two marks would commute and disconnect the parent
noncommutation diagram. Every right-hand mark consequently
acts only on $B$, and every left-hand mark only on $A$.
That is precisely a perfect split. Otherwise any shared
mark belongs to the active interval of $B$, hence has
index greater than $i$, proving the last assertion.
\end{proof}

\paragraph{Attribution and scope.}
The fracture framework and the both-imprimitive and primitive-side
routes come from Cameron--Fernandes--Leemans, \emph{The number of
string C-groups of high rank}, arXiv:2212.12723v2, Section~5,
and their earlier alternating-group source. The actual contracted
orbit forest was already established in the accepted L044 work.
Here minimal transitivity replaces the independent-group
generation step; the two commuting-factor inequalities then
give the singleton and two-point conclusions without evenness,
faithful projected factors or projected IP. The primitive
normal-subgroup argument is the acquired L044 argument, with
its exact two-point boundary now settled by this count.

The only acquired mathematical input is the full unique-crossing
and orbit-cut proof. All counting arguments are written above;
no census, rank classification, computation or unproved stability
statement is a premise. The remaining head support, split-existence
and stability obligations are separate. No whole-original,
conjecture, Lean or numerical-compression credit is asserted.
\end{document}
